\documentclass[10pt,a4paper]{amsart}
\usepackage[margin=19mm]{geometry}
\usepackage{amsmath,amssymb,mathtools}
\usepackage[scr=rsfs,cal=euler]{mathalfa}
\usepackage{xfrac}
\usepackage[unicode,hidelinks,bookmarks=false]{hyperref}
\newcommand{\ie}{i.e.\ }
\newcommand{\cf}{cf.\ }
\newcommand{\resp}{respectively\ }
\newcommand{\Subsection}{\subsection}
\providecommand{\nobreakdash}{\nobreak\hbox{-}}
\setlength{\emergencystretch}{2em}
\multlinegap0pt
\usepackage{amssymb}
\usepackage[shortlabels]{enumitem}
\usepackage{float}
\usepackage{tikz}
\usepackage{mathtools}
\newtheorem{rem}{Remark}
\newtheorem{theorem}{Theorem}
\def \Z {{\mathbb Z}}
\def \la {{\lambda}}
\title{On the Generating Graph of Finite Abelian Groups}
\author{Kavita Samant, A. Satyanarayana Reddy}
\date{Source: arXiv:2609.00102v1 (CC BY 4.0)}
\begin{document}
\maketitle

\noindent\emph{Source and licence.} On the Generating Graph of Finite Abelian Groups. Version arXiv:2609.00102v1 (CC BY 4.0), distributed by arXiv under the Creative Commons Attribution 4.0 International licence, http://creativecommons.org/licenses/by/4.0/, as recorded in the arXiv licence metadata for this exact version and shown on the abstract page https://arxiv.org/abs/2609.00102v1. Full licence notice: \texttt{licenses/arxiv-2609.00102-CC-BY-4.0.txt}.\par\smallskip

\noindent\textit{Editorial scope.} Counted unit: the article's theorem thm:lap (source0.tex lines 1301--1309). The article's complete original proof of it is delivered with it (source0.tex lines 1310--1354, one \texttt{proof} environment of 45 lines). No mathematics is written, repaired or re-derived: the statement and the proof are the article's own lines, in the article's own notation, and no retained line is substituted or re-flowed.  Retained uncounted as the source-local context the proof needs: lines 757--771.  Nothing was omitted from any retained span, so the package carries no omission record.  No retained line contains a citation, so the wrapper carries no bibliography and no bibliography entry.  No retained line contains a figure environment, an image-inclusion command or drawing code, so no image is delivered and the package declares no asset dependency.  Editorial formatting only, in addition: an AMS \texttt{amsart} preamble that restates the article's own declarations and macros used by the retained lines, namely the environments \texttt{rem}, \texttt{theorem} on separate counters, together with the standard packages those lines need.  The retained environments print in this document as Remark 1, Theorem 1 in order of appearance, whereas the article prints the same items with the numbering of its own full document.  The complete native source of the article as distributed in the arXiv e-print for this exact version is bundled unchanged as \texttt{sources/209000/source0.tex}.
\begin{rem}\label{rem:decompose}
        Let $G = S_{p_1}\times\cdots\times S_{p_k},$ where each $S_{p_i}$ is 
        a non-cyclic abelian $p_i$-group. Since $(a_1,\ldots,a_k)$ is 
        non-isolated in $\Gamma(G)$ if and only if $a_i$ is non-isolated in 
        $\Gamma(S_{p_i})$ for each $i$, we have
        $$\Gamma(G) = \otimes_{i=1}^{k}\Gamma(S_{p_i}), \qquad 
        \Delta(G) = \otimes_{i=1}^{k}\Delta(S_{p_i}).$$
        More generally, for $G = S_{p_1}\times\cdots\times S_{p_j}\times 
        \mathbb{Z}_s$ with $s = \prod\limits_{i=j+1}^{k}p_i^{u_i}$, since 
        $\widehat{\Gamma}(\mathbb{Z}_s)$ has no isolated vertices,
        $$\Gamma(G) = \left(\otimes_{i=1}^{j}\Gamma(S_{p_i})\right)
        \otimes\,\widehat{\Gamma}(\mathbb{Z}_s), \qquad \Delta(G) = 
        \left(\otimes_{i=1}^{j}\Delta(S_{p_i})\right)\otimes\,
        \widehat{\Gamma}(\mathbb{Z}_s).$$
\end{rem}

    \begin{theorem}\label{thm:lap}
Let $G$ be a group and has at least one cyclic Sylow $p$-subgroup. Then
            \[ S\left(L(\Delta(G))\right)
               = \bigcup_{i=1}^{|V(\Delta(H))|} 
                  S\left(d \widehat{D} - \lambda_i \widehat{A}\right), \]
            where $d=r_1r_2\cdots r_j,$ 
            $\widehat{D}$ is the vertex degree matrix of $\widehat{\Gamma}(\mathbb{Z}_s)$ and  
            $\widehat{A}$ is the adjacency matrix of $\widehat{\Gamma}(\mathbb{Z}_s)$.\index{Spectrum!for $\Z_n\times \Z_m$}
     \end{theorem}

    \begin{proof}
        We consider the connected component $\Delta(G)$ of $\Gamma(G)$. By 
     Remark~\ref{rem:decompose}, $\Delta(G) = \Delta(H)\otimes\widehat{
        \Gamma}(\mathbb{Z}_s)$, and so
        $$L(\Delta(G)) = D(\Delta(H))\otimes\widehat{D} - A(\Delta(H))
        \otimes\widehat{A},$$
        where $\widehat{A}$ and $\widehat{D}$ denote the adjacency and degree 
        matrices of $\widehat{\Gamma}(\mathbb{Z}_s)$ respectively. Since 
        $A(\Delta(H))$ is symmetric, there exists an orthogonal matrix $P$ 
        such that $PA(\Delta(H))P^T = \Lambda_H$, where $\Lambda_H = 
        \mathrm{diag}(\lambda_1,\ldots,\lambda_{|V(\Delta(H))|})$. Setting 
        $M = P\otimes I$, which is orthogonal since $P$ is, we have $M^{-1} 
        = M^T$, so $L(\Delta(G))$ and $M^TL(\Delta(G))M$ are similar and 
        share the same eigenvalues. Computing:
        \begin{align}\label{eq:abbound}
        M^TL(\Delta(G))M = M^T(D(\Delta(H))\otimes\widehat{D})M - 
        M^T(A(\Delta(H))\otimes\widehat{A})M.
        \end{align}
    
        Applying the mixed product property of the Kronecker product, 
        $(X\otimes Y)(U\otimes V) = XU\otimes YV$, to the right hand side 
        of Equation~\eqref{eq:abbound} with $M = P\otimes I$, we get
    \begin{align*}
        M^T(D(\Delta(H))\otimes \widehat D)M&=(P\otimes I)^T(D(\Delta(H))\otimes \widehat D)(P\otimes I)\\
        &=P^TD(\Delta(H))P\otimes \widehat D\\
        &=D(\Delta(H))\otimes \widehat D=d I\otimes \widehat D,
    \end{align*}
    \begin{align*}
      \text{and}\quad  M^T(A(\Delta(H))\otimes \widehat A)M&=(P\otimes I)^T(A(\Delta(H))\otimes \widehat A)(P\otimes I)\\
        &=P^TA(\Delta(H))P\otimes \widehat A=\Lambda_H\otimes \widehat A\\
        &=\begin{bmatrix}
            \la_1\widehat A&\bf 0&\dots &\bf 0\\
            \bf 0&\la_2\widehat A&\dots &\bf 0\\
            \vdots&\vdots&\ddots&\vdots\\
        \bf 0&\bf 0&\dots&\la_{|V(\Delta(H))|}\widehat A
        \end{bmatrix}.
    \end{align*}
    On combining the expressions, we get
    $$M^TL(\Delta(G)) M=\begin{bmatrix}
        d\widehat D-\la_1\widehat A&\bf 0&\dots &\bf 0\\
    \bf 0&d\widehat D-\la_2\widehat A&\dots &\bf 0\\
    \vdots&\vdots&\ddots&\vdots\\
    \bf 0&\bf 0&\dots&d\widehat D-\la_{|V(\Delta(H))|}\widehat A \end{bmatrix}.$$
    Thus $S(L(\Delta(G)))=S(M^TL(\Delta(G))M)=\bigcup\limits_{i=1}^{|V(\Delta(H))|}S(d \widehat D-\la_i\widehat A).$
    \end{proof}

\end{document}
